\documentclass[11pt, letterpaper]{article}
\usepackage[utf8]{inputenc}
\usepackage[T1]{fontenc}
\usepackage{amsmath}
\usepackage{amsfonts}
\usepackage{amssymb}
\usepackage{geometry}
\usepackage{verbatim}

\geometry{letterpaper, margin=1in}

\title{\textbf{Problem E: Jayden's Automaton Factory}}
\author{Time Limit: 5.0s | Memory Limit: 1024MB}
\date{}

\begin{document}

\maketitle

\section*{Description}

Having mastered the art of the manual T-Spin, Jayden has turned his attention to the theoretical limits of stacking. He is designing a machine to generate the "Perfect Loop"---an infinite sequence of pieces that keeps the board perfectly flat.

The game is played on a grid of width $W$ and infinite height.
The pieces are the standard 7 Tetrominoes (I, O, T, L, J, S, Z).
\textbf{Pieces do not rotate}. They fall from the top of the grid in a specified column $C$, settling when they collide with the floor or existing blocks. If a row becomes completely full of blocks, it clears instantly, and all blocks above it shift down by 1.

However, the sequence of pieces is not random. It is determined by a \textbf{Bag Automaton}.
The automaton is a directed graph with $N$ states (numbered $1$ to $N$). Each edge in the graph is labeled with a Piece Type $P$ and a Column Index $C$.
Transitioning along an edge $(u, v)$ with label $(P, C)$ means Jayden drops piece $P$ at column $C$.

Jayden wants to know: How many distinct paths of length $L$ exist in the automaton, starting from state 1, such that after all $L$ pieces are dropped and all line clears are processed, the board is \textbf{completely flat} (all columns have height 0)?

\section*{Input}

The first line contains three integers: $W$, $N$, and $L$.
$$ 1 \le W \le 6, \quad 1 \le N \le 50, \quad 1 \le L \le 10^{18} $$

The second line contains an integer $M$ ($1 \le M \le 200$), the number of edges in the automaton.

The next $M$ lines describe the edges. Each line contains:
$$ u \quad v \quad T \quad C $$
representing a transition from state $u$ to state $v$ by dropping piece type $T \in \{'I', 'O', 'T', 'L', 'J', 'S', 'Z'\}$ at column $C$ ($1 \le C$).
It is guaranteed that the piece fits horizontally within the width $W$ at column $C$.

\section*{Output}

Output a single integer: the number of valid paths of length $L$ modulo $10^9 + 7$.

\section*{Examples}

\subsection*{Example 1}
\textbf{Input:}
\begin{verbatim}
4 1 2
2
1 1 O 1
1 1 O 3
\end{verbatim}

\textbf{Output:}
\begin{verbatim}
2
\end{verbatim}

\textbf{Explanation:}
The board width is 4.
Path 1: Drop 'O' at col 1. Occupies cols 1-2, rows 1-2.
Path 2: Drop 'O' at col 3. Occupies cols 3-4, rows 1-2.
After both moves, rows 1 and 2 are full and clear. The board returns to flat state (0).
The reverse order (O at 3, then O at 1) is also valid. Total = 2.

\subsection*{Example 2}
\textbf{Input:}
\begin{verbatim}
3 2 1000000000000000000
2
1 2 I 1
2 1 L 1
\end{verbatim}

\textbf{Output:}
\begin{verbatim}
0
\end{verbatim}

\end{document}
